\section{Selected-Cut Consequences of the Three-Law Hypothesis}
\label{sec:temporal-realization-and-unified-field}

Part~I proves the structural calculus and continuum bridge. Part~II assumes
the explicit three-law physical hypothesis and proves its carrier-level
consequences. With the Chapter~10 three-law hypothesis fixed, the coherent
realization, fixed-bridge exactification, and endpoint-complete microscopic
structure are already in place on the selected cut data. This chapter records
the selected-cut temporal corollary, the exact visible law on the least-motion
cut, and the universal consequences that already follow before one chooses a
particular carrier realization.

The resulting structure is common to all later carrier formulations. On the selected cut,
causal order is the realized temporal interface, visible source is exchange or
returned residual from the same coherent bundle, irreversibility is
unrecovered residual, and exact visible closure rules out exogenous
constitutive completion. The sectoral laws of the later chapters are therefore
introduced as consequences of one prior selected-cut law, not as independent
starting equations.

The chapter has a narrow scope. It assumes the explicit three-law Part~II
hypothesis fixed in Chapter~10, proves the common selected-cut structure and
the universal law consequences carried by that hypothesis, and carries forward only
the Chapter~10 question of whether endpoint completion is derivable from the
first two laws.

\begin{tcolorbox}[colback=white, colframe=black!75, boxrule=0.9pt, title={Chapter 11 at a glance}]
\begin{description}[leftmargin=1.7em, itemsep=0.25em]
\item[Assumes.] The Chapter~10 three-law hypothesis on the selected cut data.
\item[Proves.] The selected-cut temporal corollary, the least-motion exact
      visible law, and the common exchange-return / irreversibility / closure
      structure used by all later carrier realizations.
\item[Open.] The only upstream law-side question carried forward is whether
      endpoint completion follows from temporal realization and terminal
      exactification.
\end{description}
\end{tcolorbox}

\subsection{Law Statement}

\begin{tcolorbox}[colback=white, colframe=black!75, boxrule=0.9pt, title={Selected-Cut Realization of Exactification}]
\begin{center}
{\large\bfseries Selected-Cut Temporal Corollary}

\medskip
\emph{On the least-motion faithful cut \(P_\ast\), physical time is the
causal realization of the temporal interface of persistent exactification.}

\medskip
\[
\begin{aligned}
t_{\mathrm{phys}}
&:= \tau_{P_\ast}(D=3,[\Theta]),\\
\mathcal L_{\mathrm{phys}}(z)
&:= P_\ast A_{\mathrm{coh}} P_\ast + z\,\CCSigmaP{P_\ast}{z}.
\end{aligned}
\]
\end{center}
\end{tcolorbox}

\begin{proposition}[Selected-cut realization form of exactification]
\label{prop:temporal-realization-law-identity}
\lean{SelectedSchur.temporal_realization_law_identity, partIV_temporal_interface_precedes_physical_time, SelectedVisibleEffectiveLaw}
Let \(P_\ast := P_{h_\ast}\) denote the least-motion faithful selected cut. The
displayed formulas are the compact statement of the present Part~II
realization theorem: they record, on one selected realization, the derived temporal
interface \((D=3,[\Theta])\) realized as physical time and the exact selected
visible effective law carried on that same cut. The displayed law is the
selected-cut realization of the earlier exactification chain exported by the
earlier structural chapters together with the continuum bridge.
\end{proposition}

(Proof in Appendix~\ref{sec:appendix-dynamics}.)

\begin{remark}[Selected-cut autonomy identity]
\label{rem:selected-cut-autonomy-identity}\lean{SelectedSchur.selected_cut_autonomy_identity, SelectedSchur.return_fromExactVisible, CoherenceCalculus.RealizedExactificationGenerator}
With the realized exactification generator \(\mathbb X\) of
Definition~\ref{def:realized-exactification-generator}, the selected-cut
realization theorem may be read as the autonomy identity
\[
P_\ast \mathbb X = \mathcal L_{\mathrm{phys}} P_\ast
\]
on the realized temporal class. In that reading, the Schur correction
\[
z\,\CCSigmaP{P_\ast}{z}
\]
is not a late auxiliary device. It is the exact internal compensator that
restores closure of the selected visible law on the faithful cut.
\end{remark}

\begin{tcolorbox}[colback=black!2, colframe=black!40, title={Operational form on the selected cut}]
\begin{description}[leftmargin=1.5em, itemsep=0.25em]
\item[Observer covariance.] Admissible characteristic observers of the same
      realized family yield horizon-preserving equivalent visible laws. Thus
      the physical observable law is well defined on the realized temporal
      class.
\item[Least-motion selection.] Among faithful characteristic cuts that support
      the same persistent visible ordering, the physically realized cut is the
      unique cut of least observer motion.
\item[Internal exchange balance.] Visible source terms are not primitive
      creations. On the physically realized cut, every admissible visible
      source term is the carrier-side boundary of an internal exchange current,
      so closed realized systems carry zero net internal creation.
\item[Residual return.] Only return acts on the physically realized cut.
      Observable forcing is the returned visible trace of residual propagated
      in the hidden/interface sector, while observable irreversibility is
      unrecovered residual.
\item[Self-sufficient closure.] On the physically realized cut, the selected
      law is exactly visible and intrinsically closed. No exogenous
      constitutive completion is admitted there; any apparent constitutive
      response is only returned residual from eliminated sectors of the same
      coherent bundle.
\end{description}
\end{tcolorbox}

\noindent
The temporal realization surface may be organized into three structural
components that recur throughout Part~II: realization-selection,
exchange-return, and intrinsic closure.

\begin{proposition}[Three-principle decomposition of the temporal realization surface]
\label{rem:five-laws-three-principles}\lean{five_law_surface_of_three_master_principles}
The temporal realization surface is equivalent to the conjunction of the
following three organizing principles. These are not additional laws beside the
same selected-cut realization surface. They are its compressed decomposition:
\begin{tcolorbox}[colback=black!2, colframe=black!35, title={Three Master Principles}]
\begin{enumerate}[label=(\roman*), leftmargin=2em, itemsep=0.2em]
\item \emph{Realization-selection}
      (\lean{RealizationSelectionPrinciple}): the selected visible law is the
      physically realized law on one coherent bundle, the realized temporal
      class is well defined, and the least-motion clause is the stability
      selector on that class.
\item \emph{Endogenous exchange-return}
      (\lean{EndogenousExchangeReturnPrinciple}): the exchange-balance and
      residual-return clauses hold, so visible source, forcing, and
      irreversibility are endogenous readouts of exchange and return on that
      same coherent bundle.
\item \emph{Self-sufficient realization}
      (\lean{SelfSufficientRealizationPrinciple}): the self-sufficient closure
      clause holds, so the selected law is already exactly visible and
      intrinsically closed on the physically realized cut, with no exogenous
      constitutive completion.
\end{enumerate}
\end{tcolorbox}
The temporal realization principle
(\lean{PhysicalRealizationPrinciple}) is the conjunction of these three
principles.
\end{proposition}

\subsection{Exact Visibility and Self-Sufficient Closure}

\begin{definition}[Temporal realization principle]
\label{def:physical-realization-principle}\lean{PhysicalRealizationPrinciple}
Fix a least-motion characteristic observer datum with selected projector
\(P_h\) and selected visible effective law \(\mathcal L_{\mathrm{vis},h}(z)\).
The realized family satisfies the \emph{temporal realization principle} on that
cut if the following three theorem-backed clauses all hold on that same datum:
\begin{enumerate}[label=(\roman*), leftmargin=2em, itemsep=0.2em]
\item the \emph{realization-selection principle} holds: the physically
      realized observable dynamics on the selected observed bundle are exactly
      \(\mathcal L_{\mathrm{vis},h}(z)\), they are read on one coherent
      bundle, the selected law persists as a visible ordered dynamics along the
      chosen horizon family, and the physically realized cut is unique among
      faithful, closure-admissible, promotion-stable characteristic observer
      data up to horizon-preserving equivalence;
\item the \emph{endogenous exchange-return principle} holds: every admissible
      visible source term on that cut is the carrier-side boundary of exchange
      with another visible sector or with the hidden complement of the same
      coherent bundle, all observable forcing on that cut is returned
      residual from hidden or interface propagation, and the irreversible
      content of the visible dynamics is exactly unrecovered residual;
\item the \emph{self-sufficient realization principle} holds on that cut.
\end{enumerate}
\end{definition}

\begin{remark}[Definition by theorem-backed shorthand]
\label{rem:temporal-realization-is-derived}\lean{SelectedSchur.derivedBridge, SelectedSchur.derivedRealization}
Definition~\ref{def:physical-realization-principle} is a shorthand, not a new
primitive bridge axiom. The selected-cut forcing theorem,
Theorem~\ref{thm:derived-temporal-interface-and-part-iii-horizon-tower-force-canonical-selected-schur-bridge},
and
Corollary~\ref{cor:derived-temporal-interface-and-part-iii-horizon-tower-force-canonical-selected-schur-realization}
below together show how this package is derived on one least-motion datum from
the earlier structural interfaces together with the continuum bridge, exact
elimination, and the Schur/HFT-2 return
layer.
\end{remark}

Corollary~\ref{cor:temporal-interface-precedes-physical-time} supplies the
ordered persistence datum. Temporal realization identifies the faithful
visible cut on which that ordered persistence is read as causal order. The
exchange-return and self-sufficient clauses then determine the meaning of
source, forcing, constitutive response, and irreversibility on that cut.

\begin{definition}[Self-sufficient realization principle]
\label{def:self-sufficient-realization-principle}\lean{SelfSufficientRealizationPrinciple}
Fix a least-motion characteristic observer datum with selected projector
\(P_h\) and selected visible effective law \(\mathcal L_{\mathrm{vis},h}(z)\).
The realized family satisfies the \emph{self-sufficient realization
principle} on that cut if:
\begin{enumerate}[label=(\roman*), leftmargin=2em, itemsep=0.2em]
\item the selected observed law is already exactly visible on the selected cut;
\item the selected observed law is intrinsically closed on that cut;
\item no exogenous constitutive completion is required on that cut;
\item any apparent constitutive term on that cut is returned residual from
      eliminated sectors of the same coherent bundle.
\end{enumerate}
\end{definition}

\begin{definition}[Exact-visible selected-cut package]
\label{def:exact-visible-selected-cut-package}\lean{SelfSufficientRealizationPrinciple, SelectedVisibleEffectiveLaw}
Fix a least-motion characteristic observer datum with selected projector
\(P_h\), coherent operator \(A_{\mathrm{coh}}\), and selected visible
effective law \(\mathcal L_{\mathrm{vis},h}(z)\). The realized family
satisfies the \emph{exact-visible selected-cut package} if:
\begin{enumerate}[label=(\roman*), leftmargin=2em, itemsep=0.2em]
\item the selected observed law on the physically realized cut is exactly
      \(\mathcal L_{\mathrm{vis},h}(z)\);
\item \(\mathcal L_{\mathrm{vis},h}(z)\) is read on that same cut by internal
      elimination of the hidden complement of the same coherent bundle, namely
      by the effective-law calculus of
      Definition~\ref{def:selected-visible-effective-law};
\item any apparent constitutive contribution on that cut is already contained
      in that internally induced hidden-complement correction and introduces no
      datum external to the same coherent bundle.
\end{enumerate}
\end{definition}

\begin{theorem}[Exact-visible selected-cut package forces self-sufficient realization]
\label{thm:exact-visible-selected-cut-package-forces-self-sufficient-realization}\lean{self_sufficient_realization_principle}
The exact-visible selected-cut package of
Definition~\ref{def:exact-visible-selected-cut-package} implies the
self-sufficient realization principle of
Definition~\ref{def:self-sufficient-realization-principle} on that same
selected cut.
\end{theorem}

(Proof in Appendix~\ref{sec:appendix-dynamics}.)

\subsection{Selection, Exchange, and Return}

The exact-visible-law condition is the mathematical form of self-sufficient
realization. Intrinsic closure and the absence of exogenous constitutive
completion are its physical reading on the time-realizing cut.

On the canonical nontrivial balanced realization exported by Part~I, the
counting bedrock singles out \(D=3\) as the unique residual interface between
the unit seed and the balanced \(D=4\) carrier
(Corollary~\ref{cor:residual-localization-on-d3}). The continuum bridge at the
end of Part~I then promotes that
unique interface to a temporal interface by combining residual localization,
scalar-recursive selected transport, and faithful continuum identification
(Theorem~\ref{thm:temporal-interface-from-d3}). The returned residual of
the earlier structural chapters together with the continuum bridge is therefore
read as a physical return field on that unique
time-bearing interface rather than as an arbitrary external source sector.

Least motion is a stability condition, not the primitive content of the law.
The selected cut is first determined by the requirement that the persistent
visible ordering of the selected law be realized as physical time. Least
motion then singles out the stable cut among faithful cuts carrying that same
ordered persistence. Here ``observer motion'' means motion in observer-data or
cut-space: the drift of the selected projector, admissible completion, and
promoted visible law across a horizon family. It does not mean physical
spacetime motion of the underlying state.

In the present development, this clause is established at the explicit finite
level: the selected cut comes from a minimizer of a finite design objective
over an admissible candidate family. The least-motion statement proved here is
therefore a finite selector theorem, not a variational theory on a continuum
of observer fields.

The exchange-balance clause does not replace closure classification. It
constrains what counts as an admissible source once the cut has been selected.
On a closed reversible carrier it forces source-free leading-order dynamics.
On coupled carriers it identifies visible sources as exchange traces with
other visible sectors or with the hidden complement.

\begin{definition}[Selected visible effective law]
\label{def:selected-visible-effective-law}\lean{SelectedVisibleEffectiveLaw}
Fix a least-motion characteristic observer datum with selected projector
\(P_h\), and let \(A_{\mathrm{coh}}\) denote the coherent operator carried by
the realized family on the ambient characteristic bundle. The
\emph{selected visible effective law} on that cut is the exact effective
operator
\[
\mathcal L_{\mathrm{vis},h}(z)
:=
P_hA_{\mathrm{coh}}P_h + z\,\CCSigmaP{P_h}{z}.
\]
Equivalently, \(\mathcal L_{\mathrm{vis},h}(z)\) is the effective operator of
Definitions~\ref{def:schur-correction} and \ref{def:effective-operator}
evaluated on the selected cut. Its \emph{selected endpoint closure class} is the
carrier-side exact closure class obtained by characteristic boundary reduction
of that same effective operator on the selected cut. Equivalently, admissible
carrier-level endpoint representatives belong to the selected endpoint closure
class precisely when they are read from the one exact visible law
\(\mathcal L_{\mathrm{vis},h}(z)\) on that selected cut and differ only by the
carrier-level continuum equivalence used later in the endpoint-rigidity
theorems. When a refinement-compatible local representative family is
exhibited, the constructive bridge theorems of
Chapter~\ref{sec:continuum-identification} verify that same class by exact
core realization or continuum forcing; they do not define the selected
endpoint class on the critical law-side route of Part~II.
\end{definition}

\begin{definition}[Scoped temporal-completion class]
\label{def:scoped-temporal-completion-class}\lean{PhysicalRealizationPrinciple}
Fix a realized family carrying the derived temporal interface exported by
Corollary~\ref{cor:temporal-interface-precedes-physical-time}. A
\emph{scoped temporal completion} of that interface is a least-motion
characteristic observer datum together with its selected visible effective law
for which the temporal realization principle holds on the selected cut. This
class is intentionally local to the present theorem surface: it ranges over
law-side completions readable on the exported structural and continuum-bridge
interfaces, not
over arbitrary continuum observer-field variational theories.
\end{definition}

\begin{proposition}[Scoped temporal completion is unique up to equivalence]
\label{prop:scoped-temporal-completion-unique}\lean{realization_selection_principle, physical_realization_principle}
Fix a realized family carrying the derived temporal interface exported by
Corollary~\ref{cor:temporal-interface-precedes-physical-time}. Assume
\((P_h,\mathcal L_{\mathrm{vis},h}(z))\) and
\((P_{h'}',\mathcal L'_{\mathrm{vis},h'}(z))\) are two members of the scoped
temporal-completion class of
Definition~\ref{def:scoped-temporal-completion-class}. Then:
\begin{enumerate}[label=(\roman*), leftmargin=2em, itemsep=0.2em]
\item the two selected cuts are horizon-preserving equivalent on the realized
      temporal class;
\item the induced observable dynamics agree up to that equivalence, so the
      selected visible law on the realized temporal class is unique modulo
      horizon-preserving equivalence;
\item the admissible reading of visible source, forcing, constitutive response,
      and irreversibility is the same on both completions: visible source is
      exchange-boundary reentry, observable forcing is returned residual, any
      apparent constitutive term is internal hidden-sector elimination, and
      observable irreversibility is unrecovered residual.
\end{enumerate}
Accordingly, on the present theorem surface there is no additional independent
law-side completion datum beyond temporal realization on the selected cut; the
remaining downstream freedom begins only at carrier visibility, carrier
generation, and carrier-level representative identification.
\end{proposition}

(Proof in Appendix~\ref{sec:appendix-dynamics}.)

\begin{corollary}[Exogenous law-side alternatives are excluded on the selected cut]
\label{cor:exogenous-law-side-alternatives-excluded}\lean{residual_return_principle, self_sufficient_realization_principle}
Fix the setting of
Definition~\ref{def:scoped-temporal-completion-class}. Any proposed visible-law
completion on the selected cut that introduces at least one of the following:
\begin{enumerate}[label=(\roman*), leftmargin=2em, itemsep=0.2em]
\item primitive visible source not read as exchange-boundary reentry;
\item observable forcing not read as returned residual from the same coherent
      bundle;
\item constitutive completion not induced by internal hidden-sector
      elimination;
\item irreversibility not read as unrecovered residual;
\end{enumerate}
is not a member of the scoped temporal-completion class.
\end{corollary}

\begin{proof}
Each prohibited item contradicts either the endogenous exchange-return clause
or the self-sufficient realization clause in
Definition~\ref{def:physical-realization-principle}. Hence such a proposed
completion fails the temporal realization principle and therefore does not lie
in the scoped temporal-completion class.
\end{proof}

\begin{definition}[Coherent-law exact elimination package]
\label{def:coherent-law-exact-elimination-package}\lean{CanonicalSelectedSchurBridge, SelectedSchur.ExactResidue}
Fix a least-motion characteristic observer datum with selected projector
\(P_h\), coherent operator \(A_{\mathrm{coh}}\), and selected visible
effective law \(\mathcal L_{\mathrm{vis},h}(z)\). The realized family
satisfies the \emph{coherent-law exact elimination package} on that cut if:
\begin{enumerate}[label=(\roman*), leftmargin=2em, itemsep=0.2em]
\item the selected observed law on the physically realized cut is exactly
      \(\mathcal L_{\mathrm{vis},h}(z)\);
\item the selected return field on that same cut is the residual return field
      of Definition~\ref{def:residual-return-field} for the coherent operator
      \(A_{\mathrm{coh}}\) and the selected projector \(P_h\);
\item every hidden-side contribution visible on that cut is read only through
      the exact Schur correction of
      Definition~\ref{def:schur-correction}, equivalently through the returned
      residual forcing of
      Proposition~\ref{prop:return-residual-is-schur}, and introduces no datum
      external to the same coherent bundle.
\end{enumerate}
\end{definition}

\begin{theorem}[Coherent-law exact elimination forces exact-visible selected-cut package]
\label{thm:coherent-law-exact-elimination-forces-exact-visible-selected-cut-package}\lean{SelectedCut.GroundedBridge.exactSelectedSchur, SelectedCut.GroundedBridge.exactVisibleSelectedLaw}
The coherent-law exact elimination package of
Definition~\ref{def:coherent-law-exact-elimination-package} implies the
exact-visible selected-cut package of
Definition~\ref{def:exact-visible-selected-cut-package} on that same selected
cut.
\end{theorem}

(Proof in Appendix~\ref{sec:appendix-dynamics}.)

\begin{remark}[Exact elimination and transferred visible law]
\label{rem:exact-elimination-and-transfer}\lean{CanonicalSelectedSchurBridge.coherentElimination, SelectedCut.GroundedBridge.exactSelectedSchur}
The viewpoint above is compatible with modern homotopy-transfer / effective
field-theory language: the visible law is not postulated independently, but is
transferred from one coherent law by exact elimination of the hidden sector
\cite{ArvanitakisHohmHullLekeu2020}. The manuscript does not use that machinery
as a black box. The proof route here is the direct Schur-return calculus
already developed in Parts~II--IV.
\end{remark}

\begin{definition}[No-auxiliary internal shadow elimination package]
\label{def:no-auxiliary-internal-shadow-elimination-package}\lean{SelfSufficientRealizationPrinciple, CanonicalSelectedSchurBridge}
Fix a least-motion characteristic observer datum with selected projector
\(P_h\), coherent operator \(A_{\mathrm{coh}}\), and selected visible
effective law \(\mathcal L_{\mathrm{vis},h}(z)\). The realized family
satisfies the \emph{no-auxiliary internal shadow elimination package} on that
cut if:
\begin{enumerate}[label=(\roman*), leftmargin=2em, itemsep=0.2em]
\item the coherent law on the selected datum decomposes into a visible equation
      and a hidden-complement/shadow equation on one and the same coherent
      bundle;
\item that shadow equation admits an exact internal solve on the selected cut
      as a unique graph over the visible datum, in the sense of
      Theorem~\ref{thm:nonlinear-shadow-elimination};
\item in the linear selected-law specialization of
      Corollary~\ref{cor:linear-shadow-specialization}, the reduced observed
      equation is exactly \(\mathcal L_{\mathrm{vis},h}(z)\), and the returned
      visible forcing is exactly the residual return / Schur correction of
      Proposition~\ref{prop:return-residual-is-schur};
\item the shadow solve depends only on the visible datum together with the
      shadow equation on that same coherent bundle, with no hidden labels,
      quotient choice, gauge slice, or exogenous constitutive datum.
\end{enumerate}
\end{definition}

\begin{theorem}[Selected visible effective law forces no-auxiliary internal shadow elimination]
\label{thm:selected-visible-effective-law-forces-no-auxiliary-shadow-elimination}\lean{physical_realization_principle, self_sufficient_realization_principle}
Fix a least-motion characteristic observer datum with selected projector
\(P_h\), coherent operator \(A_{\mathrm{coh}}\), and selected visible
effective law \(\mathcal L_{\mathrm{vis},h}(z)\) of
Definition~\ref{def:selected-visible-effective-law}. Then the same selected
cut satisfies the no-auxiliary internal shadow elimination package of
Definition~\ref{def:no-auxiliary-internal-shadow-elimination-package}.
\end{theorem}

(Proof in Appendix~\ref{sec:appendix-dynamics}.)

\begin{theorem}[No-auxiliary internal shadow elimination forces coherent-law exact elimination]
\label{thm:no-auxiliary-internal-shadow-elimination-forces-coherent-law-exact-elimination}\lean{SelectedCut.GroundedBridge.exactSelectedSchur, SelectedSchur.exactFromExactVisible}
The no-auxiliary internal shadow elimination package of
Definition~\ref{def:no-auxiliary-internal-shadow-elimination-package} implies
the coherent-law exact elimination package of
Definition~\ref{def:coherent-law-exact-elimination-package} on that same
selected cut.
\end{theorem}

(Proof in Appendix~\ref{sec:appendix-dynamics}.)

\begin{corollary}[Selected-cut forcing, no-auxiliary internal shadow elimination, and Schur/HFT-2 return force the grounded selected Schur package]
\label{cor:selected-cut-plus-no-auxiliary-shadow-plus-schur-return-forces-grounded-package}\lean{canonical_selected_schur_realization}
If, on one and the same least-motion characteristic observer datum, a realized
family satisfies:
\begin{enumerate}[label=(\alph*), leftmargin=2em, itemsep=0.2em]
\item the selected-cut forcing package of
      Definition~\ref{def:selected-cut-forcing-package};
\item the no-auxiliary internal shadow elimination package of
      Definition~\ref{def:no-auxiliary-internal-shadow-elimination-package};
\item the selected Schur/HFT-2 return package of
      Definition~\ref{def:selected-schur-hft2-return-package};
\end{enumerate}
then it satisfies the grounded selected Schur package of
Definition~\ref{def:grounded-selected-schur-package}.
\end{corollary}

\begin{proof}
By
Theorem~\ref{thm:no-auxiliary-internal-shadow-elimination-forces-coherent-law-exact-elimination},
item \emph{(b)} supplies the coherent-law exact elimination package on the same
selected cut. Together with items \emph{(a)} and \emph{(c)}, the hypotheses of
Corollary~\ref{cor:selected-cut-plus-coherent-law-exact-elimination-plus-schur-return-forces-grounded-package}
hold on one and the same datum. Therefore the grounded selected Schur package
follows.
\end{proof}

\begin{corollary}[Selected-cut forcing and Schur/HFT-2 return force the grounded selected Schur package]
\label{cor:selected-cut-plus-schur-return-forces-grounded-package}\lean{canonical_selected_schur_realization}
If, on one and the same least-motion characteristic observer datum, a realized
family satisfies:
\begin{enumerate}[label=(\alph*), leftmargin=2em, itemsep=0.2em]
\item the selected-cut forcing package of
      Definition~\ref{def:selected-cut-forcing-package};
\item the selected Schur/HFT-2 return package of
      Definition~\ref{def:selected-schur-hft2-return-package};
\end{enumerate}
then it satisfies the grounded selected Schur package of
Definition~\ref{def:grounded-selected-schur-package}.
\end{corollary}

\begin{proof}
By
Theorem~\ref{thm:selected-visible-effective-law-forces-no-auxiliary-shadow-elimination},
the same selected cut satisfies the no-auxiliary internal shadow elimination
package of
Definition~\ref{def:no-auxiliary-internal-shadow-elimination-package}. The
claim therefore follows from
Corollary~\ref{cor:selected-cut-plus-no-auxiliary-shadow-plus-schur-return-forces-grounded-package}.
\end{proof}

\begin{proposition}[The temporal realization surface as a theorem chain]
\label{prop:temporal-dissolution}\lean{barrier_D3, finite_local_type_profiles, transport_generation_yields_scalar_recursive, scalar_recursion_identifies_continuum_generator, exact_recursive_transport_yields_temporal_interface, exact_recursive_transport_forces_graph_temporal_generator, wordLocalDyadicEnvelopeRayLocalStencilInterpretation_surface, wordLocalDyadicEnvelopeRayExactCoreRealization, SelectedCut.forcingSurface, temporal_interface_precedes_physical_time, SelectedSchur.derivedBridge, SelectedSchur.derivedRealization}
Each clause of the temporal realization surface is a consequence of the
explicit theorem ladder already assembled in the manuscript's active theorem
chain:
\begin{enumerate}[label=\textup{(\roman*)}, itemsep=0.2em]
\item Corollary~\ref{cor:residual-localization-on-d3} localizes the unique
      nontrivial residual interface at \(D=3\).
\item Corollary~\ref{cor:transport-generation-yields-scalar-recursive} forces
      the selected transport channels on that interface to be scalar-recursive.
\item Theorem~\ref{thm:scalar-recursion-identifies-continuum-generator}
      constructs the unique continuum ordering generator \(\Theta\) from the
      reconstructed recursive family and proves that the same reconstruction
      maps identify the discrete scalar recursion with that continuum
      generator. When the recursive family is exact across refinement,
      Theorem~\ref{thm:exact-recursive-transport-yields-temporal-interface}
      and Corollaries~\ref{cor:exact-recursive-transport-forces-graph-temporal-generator}
      and~\ref{cor:certified-local-ray-recursion-closes-temporal-bridge}
      place that same generator on a unique temporal-interface surface without
      separate stability or summability hypotheses on the active import-free
      lane.
\item Theorem~\ref{thm:temporal-realization-forces-selected-cut-forcing-package}
      forces the unique least-motion faithful selected cut. Combined with
      Corollary~\ref{cor:temporal-interface-precedes-physical-time}, this is
      the cut on which the exported ordered persistence is read as physical
      temporal order.
\item Theorem~\ref{thm:derived-temporal-interface-and-part-iii-horizon-tower-force-canonical-selected-schur-bridge}
      and
      Corollary~\ref{cor:derived-temporal-interface-and-part-iii-horizon-tower-force-canonical-selected-schur-realization}
      then identify, on that same selected datum, exact visibility,
      endogenous exchange/return, and irreversibility as unrecovered residual.
\end{enumerate}
In particular, the temporal realization surface adds no information beyond
what this cited theorem chain already proves on the selected cut.
\end{proposition}

\begin{proof}
Items~(i)--(iii) are exactly the cited Chapter~9 bridge theorems. Item~(iv)
is the selected-cut theorem proved below from the finite least-motion
selector together with the realized-cut identification, read against the
ordered-persistence datum exported by
Corollary~\ref{cor:temporal-interface-precedes-physical-time}. Item~(v) is the derived canonical
selected Schur bridge on that same datum together with the corollary that
repackages its realization surface as the temporal realization principle.
Every clause therefore reduces to a cited theorem.
\end{proof}

\subsection{Common Grammar of Realized Dynamics}

\begin{remark}[Physical grammar on the selected cut]
\label{rem:common-part-iv-grammar}\lean{CommonPartIVGrammar}
The realized dynamics follow one common physical grammar:
\[
\begin{aligned}
\text{hidden coherent law}
&\Longrightarrow \text{persistent coherent visibility on a faithful cut}\\
&\Longrightarrow \text{physical time on the selected cut}\\
&\Longrightarrow \text{residual return field on the selected cut}\\
&\Longrightarrow \text{selected visible effective law is physical}\\
&\Longrightarrow \text{only return acts on the visible bundle}\\
&\Longrightarrow \text{characteristic endpoint reduction}\\
&\Longrightarrow \text{canonical sector generation}\\
&\Longrightarrow \text{symmetry-rigid minimal compatible closure}\\
&\Longrightarrow \text{carrier-level physical law.}
\end{aligned}
\]
This is the composition of exact projection reduction, temporal realization of
persistent exactification, stability selection, exchange-only visible sources,
constitutive emergence from eliminated sectors, and carrier-level minimal
compatible closure.

The flagship sections differ only by which carrier becomes faithfully visible
on the selected cut. In the canonical four-dimensional realization exported by
Part~II, the phase, wave, gauge, and geometric lanes are presented as
intrinsic visible-law normal forms on associated carriers of one selected
observed bundle, while the kinetic lane is presented as the effective
constitutive closure read from an intrinsic balance hierarchy on that same
bundle. Theorem~\ref{thm:common-selected-bundle-assembly} assembles those
carrier-wise realizations once they are visible on one least-motion cut. The
carrier-level step is identification of a canonical representative of the
forced physical law class.
\end{remark}

\begin{figure}[ht]
\centering
\resizebox{\textwidth}{!}{%
\begin{tikzpicture}[x=1cm,y=1cm]
  \node[ccbox, align=center, text width=3.45cm, minimum height=0.92cm] (p1) at (-4.85,5.15)
    {\textbf{Structural calculus}\\HFT grammar and elimination};
  \node[ccbox, align=center, text width=3.45cm, minimum height=0.92cm] (p2) at (0,5.15)
    {\textbf{Realized models}\\closure-balance and characteristic};
  \node[ccbox, align=center, text width=3.55cm, minimum height=0.92cm] (p3) at (4.85,5.15)
    {\textbf{Continuum bridge}\\temporal interface on the \(D=3\) residual interface};

  \draw[ccarrow, line width=0.5pt] (p1.east) -- (p2.west);
  \draw[ccarrow, line width=0.5pt] (p2.east) -- (p3.west);

  \node[ccbox, align=center, minimum width=5.6cm, minimum height=0.95cm] (realize) at (0,3.55)
    {\textbf{Temporal realization on the least-motion faithful cut}};
  \draw[ccarrow, line width=0.5pt]
    (p3.south) -- ($(realize.north east)+(-1.15,0)$);

  \node[ccbox, align=center, minimum width=3.7cm, minimum height=1.0cm] (law) at (-4.8,2.0)
    {\textbf{Selected visible law}\\is physical};
  \node[ccbox, align=center, minimum width=3.7cm, minimum height=1.0cm] (exchange) at (0,2.0)
    {\textbf{Exchange and return only}\\on the visible bundle};
  \node[ccbox, align=center, minimum width=3.7cm, minimum height=1.0cm] (closure) at (4.8,2.0)
    {\textbf{Exact visibility and}\\intrinsic closure};

  \draw[ccarrow, line width=0.5pt] (realize.south) -- (law.north);
  \draw[ccarrow, line width=0.5pt] (realize.south) -- (exchange.north);
  \draw[ccarrow, line width=0.5pt] (realize.south) -- (closure.north);

  \node[ccbox, align=center, minimum width=7.6cm, minimum height=0.78cm] (bundle) at (0,0.5)
    {\textbf{Selected observed bundle}};

  \draw[ccarrow, line width=0.5pt] (law.south) -- ++(0,-0.26) -| ($(bundle.north)+(-2.1,0)$);
  \draw[ccarrow, line width=0.5pt] (exchange.south) -- (bundle.north);
  \draw[ccarrow, line width=0.5pt] (closure.south) -- ++(0,-0.26) -| ($(bundle.north)+(2.1,0)$);

  \node[ccbox, align=center, text width=2.7cm, minimum height=1.15cm] (phase) at (-4.4,-1.05)
    {\textbf{Phase}\\primitive \(1\)\\local-state persistence};
  \node[ccbox, align=center, text width=2.7cm, minimum height=1.15cm] (wave) at (0,-1.05)
    {\textbf{Wave}\\primitive \(2\)\\propagative persistence};
  \node[ccbox, align=center, text width=2.7cm, minimum height=1.15cm] (gauge) at (4.4,-1.05)
    {\textbf{Gauge}\\primitive \(3\)\\comparative persistence};
  \node[ccbox, align=center, text width=3.05cm, minimum height=1.15cm] (geom) at (-2.2,-3.0)
    {\textbf{Geometry}\\lift of primitive \(3\)\\causal-structural response};
  \node[ccbox, align=center, text width=3.05cm, minimum height=1.15cm] (hydro) at (2.2,-3.0)
    {\textbf{Hydrodynamics}\\effective closure\\resolution-integrated balance};

  \draw[ccarrow, line width=0.5pt] ($(bundle.south)+(-2.4,0)$) -- (phase.north);
  \draw[ccarrow, line width=0.5pt] (bundle.south) -- (wave.north);
  \draw[ccarrow, line width=0.5pt] ($(bundle.south)+(2.4,0)$) -- (gauge.north);
  \draw[ccarrow, line width=0.5pt] ($(bundle.south)+(-2.2,0)$) -- (geom.north);
  \draw[ccarrow, line width=0.5pt] ($(bundle.south)+(2.2,0)$) -- (hydro.north);
\end{tikzpicture}%
}
\caption{Common physical pipeline. The structural horizon calculus, its
realized closure-balance and characteristic models, and the continuum bridge
at the end of Part~I feed the Part~II temporal realization law.
On the least-motion faithful cut the selected visible law is physical,
exchange and returned residual exhaust admissible forcing, and exact
visibility yields one selected observed bundle carrying the three primitive
direct channels, the canonical geometric lift of the primitive
\(3\)-channel, and the separate resolution-integrated hydrodynamic closure
lane.}
\label{fig:part-iv-common-pipeline}
\end{figure}

\begin{remark}[Bundle-level reading]
\label{rem:bundle-first-reading}\lean{SelectedBundle.FirstReading}
At the bundle level, one selects the faithful cut on which coherent visibility
persists, reads that ordered persistence as physical time, forms the exact
visible effective law on that cut, identifies all visible forcing as returned
residual, restricts visible sources to exchange-boundary reentry, reduces
unresolved bulk contribution to endpoint return, generates the canonical
visible carrier sectors, and classifies the surviving minimal compatible
closure on the associated visible carrier. In the canonical four-dimensional
realization of Part~II, this yields a common selected \(D=4\) observed bundle
together with a stable \(D=2\) phase sector as its smallest reversible
associated carrier.
\end{remark}

\subsection{Universal Consequences of Temporal Realization}

Before passing to carrier-level readouts, the temporal realization principle
already forces a universal physical layer. Some of these consequences are
immediate consequences of the law itself; others require the additional
variational or regularity hypotheses under which one ordinarily writes action
principles or continuous symmetry flows. They enter here as consequences of
temporal realization, not as unrelated primitive axioms.

\subsubsection{Causality}

\begin{theorem}[Temporal realization forces causal visible propagation]
\label{thm:temporal-realization-forces-causal-propagation}\lean{physical_realization_principle, residual_return_principle}
Assume the temporal realization principle on a selected cut. Then admissible
visible influence on that cut is causal in the following sense:
\begin{enumerate}[label=(\roman*), leftmargin=2em, itemsep=0.2em]
\item the visible dynamics are read on an ordered horizon family, so visible
      evolution is evaluated only along that realized order;
\item every admissible visible forcing term factors through returned residual
      propagated in the hidden or interface sector of the same coherent bundle;
\item no visible influence unsupported by that ordered propagation is
      admissible on the realized cut.
\end{enumerate}
In particular, on the physically realized cut there is no primitive acausal
visible action.
\end{theorem}

(Proof in Appendix~\ref{sec:appendix-dynamics}.)

\subsubsection{Stationarity and Action}

\begin{definition}[Stationary faithful temporal realization]
\label{def:stationary-faithful-temporal-realization}\lean{SelectedCut.GroundedBridge, observerMotionValue}
Fix a realized family and an admissible family of faithful characteristic cuts
that all carry the same persistent visible ordering. A faithful temporal
realization is \emph{stationary} if the observer-motion functional has zero
first variation at the realized cut along every admissible deformation through
that family. It is \emph{minimizing} if the realized cut attains the minimum of
that functional on the same admissible class.
\end{definition}

\begin{theorem}[Least-motion realization is stationary]
\label{thm:least-motion-realization-is-stationary}\lean{SelectedCut.groundedBridge, observerMotionValue_minimal}
Assume the temporal realization principle on a realized family, and assume the
admissible faithful cuts carrying the same persistent visible ordering form a
differentiable family on which the observer-motion functional is differentiable.
Then the physically realized least-motion cut is stationary in the sense of
Definition~\ref{def:stationary-faithful-temporal-realization}. If the least-
motion cut is unique, it is the unique minimizing stationary faithful
realization.
\end{theorem}

\begin{proof}
By the temporal realization principle, the physically realized cut is selected
as the unique least-motion cut among faithful cuts carrying the same ordered
persistent visible law. A differentiable minimizer of a differentiable
functional has vanishing first variation. Hence the realized cut is stationary.
If the least-motion cut is unique, it is also the unique minimizing stationary
realization on that admissible family.
\end{proof}

\subsubsection{Symmetry and Conservation}

\begin{corollary}[Stationary action on reversible realized sectors]
\label{cor:stationary-action-on-reversible-realized-sectors}\lean{PhaseLane.leastMotionObserver, WaveLane.leastMotionObserver, GaugeLane.leastMotionObserver, GeometricLane.leastMotionObserver}
Assume the setting of
Theorem~\ref{thm:least-motion-realization-is-stationary}. Suppose in addition
that on a reversible visible carrier the observer-motion functional is
represented on admissible histories by a differentiable action functional
\[
S[\phi].
\]
Then the physically realized history on that carrier satisfies
\[
\delta S = 0.
\]
If the action is local and regular, the corresponding Euler--Lagrange
equations are exactly the reversible carrier law on that realized sector.
\end{corollary}

\begin{proof}
By Theorem~\ref{thm:least-motion-realization-is-stationary}, the realized
history is stationary for the observer-motion functional. If that functional is
represented by the action \(S[\phi]\) on the reversible carrier, then the
first variation of \(S\) vanishes at the realized history. Under the standard
local regularity assumptions, vanishing of the first variation yields the
Euler--Lagrange equations. Thus stationary action is a reversible-sector form
of temporal realization rather than an extra primitive postulate.
\end{proof}

\begin{definition}[Temporal stabilizer symmetry]
\label{def:temporal-stabilizer-symmetry}\lean{RealizationSelectionPrinciple, EndogenousExchangeReturnPrinciple}
Fix a realized family satisfying the temporal realization principle on a
selected cut. A \emph{temporal stabilizer symmetry} is an admissible
transformation of the realized data that preserves:
\begin{enumerate}[label=(\roman*), leftmargin=2em, itemsep=0.2em]
\item the realized temporal class of faithful cuts up to horizon-preserving
      equivalence;
\item the selected visible effective law on that class;
\item the exchange / residual-return bookkeeping on that class;
\item the inherited endpoint closure class on the visible carrier under
      consideration.
\end{enumerate}
\end{definition}

\begin{theorem}[Physical symmetries are stabilizers of temporal realization]
\label{thm:physical-symmetries-are-temporal-stabilizers}\lean{realization_selection_principle}
Assume the temporal realization principle on a realized family. Then the
physical symmetries of the visible dynamics are exactly the admissible
transformations that stabilize the realized temporal structure in the sense of
Definition~\ref{def:temporal-stabilizer-symmetry}. In particular, symmetry is
not an extra declaration beside the law: it is the stabilizer of the
temporally realized coherent dynamics.
\end{theorem}

(Proof in Appendix~\ref{sec:appendix-dynamics}.)

\begin{theorem}[Noether-type bridge on temporally realized reversible sectors]
\label{thm:noether-bridge-on-temporally-realized-reversible-sectors}\lean{realization_selection_principle}
Assume the hypotheses of
Corollary~\ref{cor:stationary-action-on-reversible-realized-sectors}. Suppose
in addition that a one-parameter family of temporal stabilizer symmetries acts
smoothly on the admissible histories and preserves the action functional
\(S\). Then the realized reversible sector carries a conserved visible current
\(J\) satisfying
\[
\partial J = 0
\]
for the carrier-side boundary / divergence / codifferential operator
appropriate to that sector. Equivalently, the corresponding Noether charge is
constant along the realized temporal evolution.
\end{theorem}

(Proof in Appendix~\ref{sec:appendix-dynamics}.)

\subsubsection{Irreversibility}

\begin{theorem}[Temporal realization forbids primitive internal creation]
\label{thm:temporal-realization-forbids-primitive-internal-creation}\lean{endogenous_exchange_return_principle, coupled_exchange_balance}
Assume the temporal realization principle on a finite coupled family of visible
carriers on one selected bundle. Then every admissible visible source term is
an exchange term with another visible sector or with the hidden complement of
the same coherent bundle. In particular, a closed temporally realized system
has no net internal creation on the visible side.
\end{theorem}

\begin{proof}
This is exactly the exchange-balance clause of the temporal realization
principle. By that clause, visible sources are not primitive creations but
boundary readouts of internal exchange on the same coherent bundle. Therefore a
closed coupled visible system cannot carry net internal creation. The explicit
finite coupled cancellation formula is written below in
Theorem~\ref{thm:coupled-exchange-balance}.
\end{proof}

\begin{theorem}[Unrecovered residual orients physical time]
\label{thm:unrecovered-residual-orients-physical-time}\lean{SelectedCut.effectiveThermodynamicClosure, residual_return_principle}
Assume the temporal realization principle on a selected cut, and assume the
selected cut carries a quantitative unrecovered-residual functional whose
monotone part measures visible irreversibility. Then:
\begin{enumerate}[label=(\roman*), leftmargin=2em, itemsep=0.2em]
\item irreversible visible evolution is oriented by the monotone accumulation
      of unrecovered residual;
\item reversible realized sectors are exactly those on which no monotone
      unrecovered residual accumulates at the visible level;
\item the thermodynamic arrow of time is therefore not an extra primitive law,
      but the macroscopic reading of residual nonrecovery on the realized cut.
\end{enumerate}
\end{theorem}

\begin{proof}
By the temporal realization principle, observable irreversibility is exactly
unrecovered residual on the realized cut. If the selected cut carries a
quantitative unrecovered-residual functional whose monotone part measures that
irreversibility, then the direction of monotone increase determines the
orientation of irreversible visible evolution. Conversely, where no monotone
unrecovered residual accumulates, the visible law is reversible at the
realized level. Thus the arrow of time is read from residual nonrecovery on the
same cut on which physical time is realized.
\end{proof}

\subsubsection{Standard Structural Consequences}

The preceding results are intended to recover the standard structural backbone
of physics before any carrier-specific equation is recognized. To keep the
scope exact, the next proposition records only what has already been proved at
this universal level and distinguishes it from the later carrier-level
identification problem.

\begin{proposition}[Scoped recovery of standard physical principles]
\label{prop:scoped-recovery-of-standard-principles}\lean{physical_realization_principle, SelectedCut.effectiveThermodynamicClosure, coupled_exchange_balance}
Assume the temporal realization principle on a selected cut, together with the
additional hypotheses stated in the corresponding results above when required.
Then the following standard structural principles are recovered as
consequences of temporal realization:
\begin{enumerate}[label=(\roman*), leftmargin=2em, itemsep=0.2em]
\item \emph{Causal propagation / locality:}
      Theorem~\ref{thm:temporal-realization-forces-causal-propagation} shows
      that visible influence on the realized cut is causal and that no
      primitive acausal visible action is admissible.
\item \emph{Stationary realization / action principle:}
      Theorem~\ref{thm:least-motion-realization-is-stationary} and
      Corollary~\ref{cor:stationary-action-on-reversible-realized-sectors}
      show that least-motion temporal realization becomes stationary
      realization, and on reversible sectors with an action representation this
      yields the usual stationary-action form \(\delta S=0\).
\item \emph{Symmetry principle:}
      Theorem~\ref{thm:physical-symmetries-are-temporal-stabilizers} identifies
      physical symmetries with stabilizers of the temporally realized coherent
      dynamics.
\item \emph{Noether bridge:}
      Theorem~\ref{thm:noether-bridge-on-temporally-realized-reversible-sectors}
      shows that continuous temporal stabilizers of a stationary reversible
      sector yield conserved visible currents.
\item \emph{Conservation / exchange balance:}
      Theorem~\ref{thm:temporal-realization-forbids-primitive-internal-creation}
      shows that closed temporally realized systems carry no primitive visible
      creation, and that admissible visible sources are internal exchange or
      returned residual.
\item \emph{Arrow of time / irreversibility:}
      Theorem~\ref{thm:unrecovered-residual-orients-physical-time} shows that
      the macroscopic arrow of time is the monotone reading of unrecovered
      residual on the realized cut.
\end{enumerate}
These recovered principles are universal and carrier-generic. They do not yet
identify the canonical representatives of the phase, wave, kinetic, gauge, or
geometric endpoint classes. That further identification is the distinct
carrier-level task taken up below.
\end{proposition}

\begin{proof}
Each item is exactly the cited theorem or corollary. The final sentence records
the distinction between universal structural recovery and carrier-level
classification. The former has been proved in the present subsection; the
latter is deferred to the sector-generation and flagship layers that follow.
\end{proof}

\begin{remark}[Scope of universal recovery]
\label{rem:what-universal-recovery-does-not-claim}\lean{physical_realization_principle, flagship_law_completion}
The proposition above isolates a carrier-independent structural layer. It
recovers causal propagation, stationary realization, symmetry, Noether-type
conservation, exchange-balanced source admissibility, and irreversibility
before any carrier-specific endpoint equation is identified. The recognizable
Schr\"odinger, wave/Klein--Gordon, Maxwell, Einstein, and effective
hydrodynamic laws enter only after the carrier-level closure classes are
generated and their canonical representatives are identified.
\end{remark}


% --- Selected-cut framework and physical dynamics (continuation) ---

